% ============================================================
% Subsection: Adaptive Strategies
% XCEPT Article 1 — Section 4
% Place this immediately after the Timescale Expansion subsection
% ============================================================

\subsection{Adaptive Strategies}
\label{sec:adaptive}

The timescale analysis of Section~\ref{sec:timescale} reveals that different
modes become operative at different resolutions and at different moments of the
trajectory.  A natural question is whether an actor can exploit this temporal
structure deliberately: \emph{enable the right modes at the right times, allow
sufficient iterations for each mode to accumulate enough firing to be effective,
and switch modes when the state signals that the current regime has been
exhausted.}  We call such a plan an \emph{adaptive strategy}.

\paragraph{Formalisation.}
An adaptive strategy for the peace actor partitions the simulation horizon into
consecutive \emph{regimes} $\{B_1, B_2, \ldots\}$, each associated with a
subset of active reactions.  At each step the actor evaluates a \emph{regime
indicator}---a scalar summary of the current state---and switches to a new
regime when the indicator crosses a threshold.  Within each regime, the
within-pool allocation fractions $(\alpha, \gamma)$ are set so that only the
targeted mode is funded: $\alpha = 0$ suppresses institution-building ($c_3$),
while $\gamma = 0$ suppresses armed-group suppression ($c_4$).

In the example that follows we implement the adaptive strategy as a
\emph{Monte-Carlo model-predictive-control} (MC-MPC) rule: at each step $t$,
$n_{\mathrm{s}} = 30$ candidate $(\alpha, \gamma)$ pairs are sampled uniformly
from $[0,1]^2$; for each candidate the deterministic Michaelis-Menten kinetics
are rolled forward for $H = 50$ steps and scored by the peace advantage
$\pi(\mathbf{x}) = (I + T + R_p) - (G + A + R_c)$; the candidate with the
highest score is then applied in the actual \emph{stochastic} step.  The
stochastic budget-allocation dynamics described in
Appendix~A draw random fractions of the available stocks $A$ and $R_c$
(conflict actor) and $R_p$ and $I$ (peace actor) as investment budgets at
each iteration; the total quantities consumed from $R_p$ and $R_c$ across the
simulation therefore serve as a natural \emph{cost of implementation}:
$R_c$ tracks the material resources the conflict actor spends, and $R_p$ tracks
the peace actor's expenditure.

\paragraph{Scenario: no-policy fails, fixed policy fails, adaptive succeeds.}
Figure~\ref{fig:three_strategies} compares three strategies starting from the
same initial state $\mathbf{x}_0 = (G,A,I,T,R_c,R_p) = (4,2,2,5,4,6)$.

\begin{itemize}
\item \textbf{No policy} $(\alpha = \gamma = 0)$: the peace actor makes no
      investment.  With no suppression and no institution-building, the conflict
      sustaining mode $C2$ operates unchallenged.  Armed groups grow to the
      carrying capacity ($A \to 15$) and institutions collapse ($I \to 0$).

\item \textbf{Fixed policy} $(\alpha = 0,\; \gamma = 0.8)$: the peace actor
      commits $80\%$ of its $I$-budget to suppression ($c_4$) and none to
      institution-building ($c_3$).  Initially this activates the suppression
      mode $P4$, reducing $A$.  However, because $c_3 = 0$, mode $P2/P3$ is
      permanently disabled.  Institutions can only be consumed by the fighting
      reaction $c_4$ (which is catalysed by $I$), so $I$ erodes monotonically.
      Once $I \to 0$, the suppression mode $P4$ itself collapses---$c_4$
      requires $I$ as a co-reactant---and $A$ recovers.  The fixed policy
      therefore fails at the same long-run outcome as no-policy, merely
      delaying collapse.

\item \textbf{MC adaptive}: the lookahead consistently selects $\alpha \approx
      0.7$--$0.9$ alongside a moderate $\gamma$, discovering that institution
      building must accompany suppression to prevent $I$ from being eroded.
      After $\approx 50$ steps the system crosses into the regime where $A < 1$
      and $I > 10$; thereafter the mode analysis (Figure~\ref{fig:mc_modes})
      confirms that only the peace-sustaining mode $P3$ remains active at
      longer timescales, consolidating the peace outcome.
\end{itemize}

\begin{figure}[h]
  \centering
  \includegraphics[width=\linewidth]{figures/5k_three_strategies.png}
  \caption{Three-strategy comparison under stochastic budget-allocation kinetics
    (300 steps, $\mathbf{x}_0 = (4,2,2,5,4,6)$).  No policy: conflict wins.
    Fixed policy ($\alpha{=}0,\,\gamma{=}0.8$): institutions collapse once $I$
    is eroded, and armed groups recover.  MC adaptive: sustained peace is
    reached, with institutions and territory growing to carrying capacity.}
  \label{fig:three_strategies}
\end{figure}

\paragraph{Cost analysis.}
Figure~\ref{fig:policy_cost} shows the cumulative $c_3$ (institution-building)
and $c_4$ (security) fluxes for each strategy.  The fixed policy spends
$\approx 6.5$ units on security ($c_4$) and zero on institution-building
($c_3$).  The MC adaptive policy spends $\approx 14.2$ units on $c_3$ and
$\approx 6.1$ units on $c_4$---a larger total, but it achieves peace.
Crucially, the \emph{fixed policy's resource expenditure stops early}: once $I
\to 0$ the suppression mode is inoperative and no further $c_4$ investment can
be absorbed.  The adaptive strategy therefore does not spend much more than the
fixed policy in absolute terms; the difference is that it allocates resources
to both modes in the right temporal order, allowing each to complete before
switching.

\begin{figure}[h]
  \centering
  \includegraphics[width=\linewidth]{figures/5m_policy_cost.png}
  \caption{Cumulative cost split into $c_3$ (institution-building, left) and
    $c_4$ (security, right) for the three strategies.  Dashed lines show
    deterministic single-lever extremes (pure security / pure growth) for
    reference.  The adaptive strategy invests more in $c_3$ than the fixed
    policy, but achieves a qualitatively different outcome at a similar total
    resource cost.}
  \label{fig:policy_cost}
\end{figure}

\paragraph{Mode activation over time.}
Figure~\ref{fig:mc_modes} shows the timescale-resolved mode analysis for the
three strategies.  In the \emph{no-policy} run the conflict markers ($G,A,R_c$)
remain in overproduction or challenge mode at all timescales throughout the
simulation, while peace markers are persistently in problem mode.  In the
\emph{fixed-policy} run the peace markers show a brief window of challenge or
overproduction at short timescales (when suppression first fires), but revert
to problem mode at longer timescales as $I$ depletes.  The conflict markers
partially enter problem mode early (suppression is working), yet recover once
$I$ can no longer sustain $c_4$.  In the \emph{MC adaptive} run the mode
structure is qualitatively different: the peace markers transition through
challenge and into overproduction at progressively longer timescales as the
run progresses, while the conflict markers move from challenge to problem mode
and stay there.  This transition is not instantaneous---it unfolds over roughly
50 steps before the mode locks in---illustrating the core principle of the
adaptive strategy: one must \emph{wait} for a mode to accumulate sufficient
activation before testing whether to switch.

\begin{figure}[h]
  \centering
  \includegraphics[width=0.49\linewidth]{figures/5n_np_mode_conflict.png}%
  \includegraphics[width=0.49\linewidth]{figures/5n_np_mode_peace.png}\\[4pt]
  \includegraphics[width=0.49\linewidth]{figures/5n_pp_mode_conflict.png}%
  \includegraphics[width=0.49\linewidth]{figures/5n_pp_mode_peace.png}\\[4pt]
  \includegraphics[width=0.49\linewidth]{figures/5n_mc_mode_conflict.png}%
  \includegraphics[width=0.49\linewidth]{figures/5n_mc_mode_peace.png}
  \caption{Timescale-resolved mode classification (overproduction/challenge/problem)
    for conflict species $\{G,A,R_c\}$ (left column) and peace species
    $\{I,T,R_p\}$ (right column), across the three strategies: no-policy (top
    row), fixed policy (middle row), MC adaptive (bottom row).  The adaptive run
    is the only one in which peace markers transition into persistent
    overproduction at longer timescales while conflict markers settle into problem
    mode.}
  \label{fig:mc_modes}
\end{figure}

\paragraph{Design principle.}
The example illustrates a general design template for adaptive strategies in
reaction-network conflict models:

\begin{enumerate}
  \item \textbf{Identify operative modes.} Use the timescale expansion
        (Section~\ref{sec:timescale}) to determine which modes are structurally
        present in each phase of the dynamics.

  \item \textbf{Enable modes sequentially.} Activate the mode that addresses the
        current bottleneck (suppression while $A$ is large; institution-building
        while $I$ is low), and commit enough resources for the mode to fire a
        substantial number of times before re-evaluating.

  \item \textbf{Monitor the regime indicator.} Track a compact state summary
        (here, $\pi(\mathbf{x}) = (I+T+R_p) - (G+A+R_c)$, or the enabling
        degrees $\pi_{P4} = \min(I,A)$ and $\alpha_{C2} = \min(G,R_c)$) to
        detect when the current mode has been exhausted or when a new threat
        has emerged.

  \item \textbf{Switch conservatively.} Hysteresis margins around the switching
        threshold prevent chattering between regimes when the state fluctuates
        near the boundary, ensuring each mode accumulates enough activation to
        have a measurable stoichiometric effect.
\end{enumerate}

This template is not specific to the MC-MPC implementation used here; any
mechanism that gates reactions based on the current state and waits
sufficiently long before switching implements the same logic.  The reaction
network formalism makes the template explicit: the stoichiometric matrix $S$,
the enabling degrees, and the mode projections are all computable directly from
the network specification, so the design of an adaptive strategy can be guided
by structural analysis rather than trial and error.
